\chapter{System Design}

\section{Overview}

This chapter describes how the system is structured. It covers the overall architecture, the
feature extraction pipeline, the classification process, and the ambiguity detection mechanism.
It also explains the design decisions and the reasoning behind them.

\section{System Architecture Overview}

The system operates in two main phases: the training phase and the inference phase. The
training phase is performed once to build and save the model, while the inference phase is
executed for each new input provided by the user.

\textbf{Figure~\ref{fig:arch}} illustrates the overall system architecture.

%% ---- Figure 4.1: System Architecture ----
\begin{figure}[H]
\centering
\begin{tikzpicture}[
  font=\small,
  box/.style={draw, rectangle, rounded corners=3pt, thick, align=center, minimum width=3.5cm, minimum height=1.0cm, fill=blue!8},
  widebox/.style={draw, rectangle, rounded corners=3pt, thick, align=center, minimum width=5.5cm, minimum height=1.0cm, fill=teal!10},
  arr/.style={-{Stealth[length=6pt]}, thick}
]

  % Nodes - Increased vertical and horizontal spacing to extend arrows
  \node[box, fill=green!12] (input) at (0,0) {\textbf{Input Text}};
  
  \node[box] (embed) at (0,-2.0) {\textbf{Sentence Embedding}\\[2pt](all-mpnet-base-v2)};
  
  \node[box] (tfidf) at (-4.2,-4.5) {\textbf{TF-IDF Features}\\[2pt](3000-dimensional)};
  
  \node[box] (knn)   at (4.2,-4.0) {\textbf{K-Nearest}\\[2pt]Neighbours (K=24)};
  \node[box] (vr)    at (4.2,-6.0) {\textbf{Vietoris-Rips}\\[2pt]Persistence (H0, H1)};
  \node[box] (pl)    at (4.2,-8.0) {\textbf{Persistence}\\[2pt]Landscape};
  
  \node[widebox] (fusion) at (0,-10.5) {\textbf{Feature Fusion}\\[2pt]\textbf{(Embedding $|$ TF-IDF $|$ Topology)}};
  
  \node[box, fill=purple!10] (xgb) at (0,-12.5) {\textbf{XGBoost Classifiers}\\[2pt]\textbf{(One-vs-Rest)}};
  
  \node[box, fill=orange!12] (ambig) at (0,-14.5) {\textbf{Ambiguity Detection}\\[2pt]\textbf{(Loop Count + Confidence Gap)}};
  
  \node[box, fill=red!8] (output) at (0,-16.5) {\textbf{Final Output}\\[2pt]\textbf{(POS / NEG / NEU / AMBIG)}};

  % Arrows - Exact routing as per reference image
  \draw[arr] (input) -- (embed);
  
  \draw[arr] (embed.west) -| (tfidf.north);
  \draw[arr] (embed.east) -| (knn.north);
  \draw[arr] (embed.south) -- (fusion.north);
  
  \draw[arr] (knn) -- (vr);
  \draw[arr] (vr) -- (pl);
  
  % Routes into the sides of Fusion box
  \draw[arr] (tfidf.south) |- (fusion.west);
  \draw[arr] (pl.south)    |- (fusion.east);
  
  \draw[arr] (fusion) -- (xgb);
  \draw[arr] (xgb) -- (ambig);
  \draw[arr] (ambig) -- (output);

\end{tikzpicture}
\caption{System Architecture for Emotion Detection with Ambiguity Detection}
\label{fig:arch}
\end{figure}

As shown in \textbf{Figure~\ref{fig:arch}}, the system processes the input text through three distinct
parallel tracks after an initial embedding step. The dense embeddings, TF-IDF representations,
and topological persistence landscapes are fused into a single comprehensive feature vector.
This vector is then passed to the XGBoost classifiers and finally evaluated by the ambiguity
detection layer to yield the final predicted label.

\section{Feature Extraction Design}

Three types of features are extracted for each input. Each type captures a different aspect of the
text.

\subsection{Semantic Embedding}

Text inputs are converted into dense vector representations using a pretrained sentence
embedding model. These embeddings capture semantic meaning in a high-dimensional space.

To quantify similarity between embeddings, we use the Euclidean distance:
\begin{equation}
d(x_i, x_j) = \|x_i - x_j\|_2
\end{equation}

This distance metric is later used to construct topological structures over the embedding space.

\subsection{TF-IDF Feature}

In addition to semantic embeddings, textual features are extracted using the Term
Frequency--Inverse Document Frequency (TF-IDF) representation.

\begin{equation}
w(t, d) = tf(t, d) \cdot \log\left(\frac{N}{df(t)}\right)
\end{equation}

where $tf(t, d)$ is the term frequency of term $t$ in document $d$, $df(t)$ is the document
frequency, and $N$ is the total number of documents.

This representation captures the importance of words within the corpus.

\subsection{Topological Features}

To capture the structural complexity of the embedding space, we employ techniques from
Topological Data Analysis (TDA). This approach analyzes the shape of data by studying how
points connect across multiple scales.

\paragraph{Vietoris--Rips Complex.}
A simplicial complex is constructed by connecting points whose pairwise distance is below a
threshold $\epsilon$:
\begin{equation}
\text{Edge exists if } d(x_i, x_j) \leq \epsilon
\end{equation}

\paragraph{Persistent Homology.}
Topological features such as connected components and loops are tracked across different
scales. Each feature is characterized by its birth and death:
\begin{equation}
\text{Persistence} = d_{\text{death}} - d_{\text{birth}}
\end{equation}

Features with high persistence represent significant structure, while those with low persistence
are considered noise.

\paragraph{Loop Features ($H_1$).}
In this work, we focus on one-dimensional homology ($H_1$), which corresponds to loops in the
data:
\begin{equation}
H_1 = \text{Number of 1-dimensional holes (loops)}
\end{equation}

A higher number of loops indicates greater structural complexity in the embedding space.

\paragraph{Persistence Landscape.}
To transform topological information into a vector representation, persistence diagrams are
converted into persistence landscapes:
\begin{equation}
\lambda_k(t) = k\text{-th largest value of } \left\{ \min(t - b_i,\; d_i - t) \right\}_+
\end{equation}

where $(b_i, d_i)$ denote the birth and death of the $i$-th topological feature, and
$(\cdot)_+$ denotes the positive part.

These topological features are later combined with embedding and TF-IDF features for
classification and ambiguity detection.

%% ---- Figure 4.2: Feature Extraction Pipeline ----
\begin{figure}[H]
\centering
\begin{tikzpicture}[
  font=\small,
  box/.style={draw, rectangle, thick, align=center, minimum width=3.5cm, minimum height=1.5cm, fill=blue!10},
  widebox/.style={draw, rectangle, rounded corners=5pt, thick, align=center, minimum width=10cm, minimum height=0.8cm, fill=teal!12},
  arr/.style={-{Stealth[length=5pt]}, thick}
]

  \node[widebox] (raw) at (0,0) {Raw Input Text};

  \node[box] (st)    at (-3.8,-3.2) {SentenceTransformer\\[2pt](768-dimensional)};
  \node[box] (tfidf) at (0,-3.2)    {TF-IDF Vectorizer\\[2pt](3000-dimensional)};
  \node[box] (topo)  at (3.8,-3.2)  {Local VR Topology\\[2pt](30-dimensional\\Landscape)};

  \node[widebox] (concat) at (0,-6.4) {Concatenated Feature Vector (3798 dimensions)};

  \draw[arr] (raw.south -| st.north) -- (st.north);
  \draw[arr] (raw.south) -- (tfidf.north);
  \draw[arr] (raw.south -| topo.north) -- (topo.north);
  
  \draw[arr] (st.south) -- (concat.north -| st.south);
  \draw[arr] (tfidf.south) -- (concat.north);
  \draw[arr] (topo.south) -- (concat.north -| topo.south);

\end{tikzpicture}
\caption{Feature Extraction Pipeline}
\label{fig:feat}
\end{figure}

\textbf{Figure~\ref{fig:feat}} provides a detailed view of the feature extraction pipeline. The raw text
is simultaneously processed by the embedding model to capture semantic meaning, the TF-IDF
vectorizer to capture keyword occurrences, and the topological data analysis module to capture
geometric complexity. These three distinct feature sets are concatenated into a unified
3798-dimensional vector utilized by the downstream classifier.

\section{Design Justification}

\subsection{Why Topology Instead of Only Confidence}

Classifier confidence alone is not reliable for detecting ambiguity. A model may produce
high confidence even for unclear inputs.

Topological features provide a geometric perspective. If neighbouring points belong to different
emotion clusters, the structure becomes complex (more loops), indicating ambiguity.

\subsection{Why 24 Nearest Neighbours}

The number of neighbours was chosen empirically:
\begin{itemize}
    \item Too few neighbours $\rightarrow$ insufficient structure
    \item Too many neighbours $\rightarrow$ loss of locality
\end{itemize}
A value of 24 provides stable and meaningful topological representations.

\subsection{Why SMOTE Was Not Used}

SMOTE generates synthetic samples by interpolation. This alters the geometry of the data.
Since TDA depends on the real geometric structure, adding synthetic points would distort
topological features and reduce reliability.

Instead, class imbalance is handled through threshold tuning.

\section{Classification Design}

Three binary XGBoost classifiers are trained using a one-vs-rest approach:
\begin{itemize}
    \item Positive classifier
    \item Negative classifier
    \item Neutral classifier
\end{itemize}

Each classifier outputs a probability score. The highest probability determines the predicted
class, subject to ambiguity checks.

XGBoost is chosen because:
\begin{itemize}
    \item It handles mixed feature types well
    \item It is robust to feature scaling differences
    \item It performs efficiently on structured data
\end{itemize}

\section{Ambiguity Detection Design}

Ambiguity detection is based on:
\begin{itemize}
    \item Loop count (topological complexity)
    \item Confidence gap (difference between top two probabilities)
\end{itemize}

The decision rules are:
\begin{enumerate}
    \item Loop count $\geq 14 \Rightarrow$ Ambiguous
    \item Loop count $\geq 10$ and confidence gap $< 0.30 \Rightarrow$ Ambiguous
    \item Confidence gap $< 0.10 \Rightarrow$ Ambiguous
    \item Otherwise $\Rightarrow$ Predicted class
\end{enumerate}

%% ---- Figure 4.3: Ambiguity Detection Flowchart ----
\begin{figure}[H]
\centering
\begin{tikzpicture}[
  font=\small,
  decision/.style={draw, diamond, thick, align=center, aspect=1.5, minimum width=2.5cm, fill=yellow!15},
  block/.style={draw, rectangle, thick, align=center, minimum width=3cm, minimum height=0.8cm, fill=red!10},
  start/.style={draw, rectangle, thick, align=center, rounded corners=10pt, minimum width=3cm, minimum height=0.6cm, fill=green!12},
  arr/.style={-{Stealth[length=5pt]}, thick}
]

  \node[start] (s) at (0,0) {Start Ambiguity Check};
  
  \node[decision] (d1) at (0,-2.5) {Loops $\geq 14$?};
  \node[block] (r1) at (5,-2.5) {Return AMBIGUOUS};
  
  \node[decision] (d2) at (0,-6) {Loops\\$\geq 10$ AND\\Gap $< 0.30$?};
  \node[block] (r2) at (5,-6) {Return AMBIGUOUS};
  
  \node[decision] (d3) at (0,-9.5) {Gap $< 0.10$?};
  \node[block] (r3) at (5,-9.5) {Return AMBIGUOUS};
  
  \node[block, fill=green!12] (r4) at (0,-12.5) {Return Highest\\Probability Class};

  \draw[arr] (s) -- (d1);
  
  \draw[arr] (d1) -- node[above=5pt, font=\scriptsize\bfseries, pos=0.4]{Yes} (r1);
  \draw[arr] (d1) -- node[left=5pt, font=\scriptsize\bfseries]{No} (d2);
  
  \draw[arr] (d2) -- node[above=5pt, font=\scriptsize\bfseries, pos=0.4]{Yes} (r2);
  \draw[arr] (d2) -- node[left=5pt, font=\scriptsize\bfseries]{No} (d3);
  
  \draw[arr] (d3) -- node[above=5pt, font=\scriptsize\bfseries, pos=0.4]{Yes} (r3);
  \draw[arr] (d3) -- node[left=5pt, font=\scriptsize\bfseries]{No} (r4);

\end{tikzpicture}
\caption{Ambiguity Detection Activity Flowchart}
\label{fig:ambig}
\end{figure}

As depicted in \textbf{Figure~\ref{fig:ambig}}, the ambiguity detection process follows a strict
sequential decision tree. It first checks for extreme geometric complexity via the loop count
before evaluating moderate complexity combined with the classifier's internal confidence gap.
If none of the ambiguity conditions are triggered, the system defaults to outputting the class
with the highest probability.

\section{Data Flow Modeling}

To visualize how data moves through the system, we employ Data Flow Diagrams (DFD) at
multiple levels. This approach allows for both a high-level overview of system interactions and
a detailed view of internal data transformations.

\subsection{Level 0: Context Diagram}

The Context Diagram (Level 0 DFD) represents the entire system as a single process and
illustrates its relationship with external entities. In this project, the primary external entity is
the User.

%% ---- Figure 4.4: Level 0 Context Diagram ----
\begin{figure}[H]
\centering
\begin{tikzpicture}[
  font=\small,
  box/.style={draw, rectangle, rounded corners=5pt, thick, align=center,
              minimum width=3.2cm, minimum height=1.6cm, fill=green!10},
  circ/.style={draw, circle, thick, align=center, minimum size=3.8cm, fill=blue!8},
  arr/.style={-{Stealth[length=6pt]}, thick}
]

% Increased horizontal distance to extend arrows
\node[box] (user) at (0,0) {%
  \textbf{User / External}\\[2pt]
  \textbf{Application}
};

\node[circ] (sys) at (8.5,0) {%
  \textbf{0.}\\[4pt]
  \textbf{Emotion Detection}\\[2pt]
  \textbf{\& Ambiguity}\\[2pt]
  \textbf{Analysis System}
};

% Top arrow: User -> System (tilted)
% Using yshift to move label further from the line
\draw[arr] (user.30) -- node[above=6pt, sloped, pos=0.5]{\small Raw Text Input} (sys.155);

% Bottom arrow: System -> User (tilted)
% Using yshift to move label further from the line
\draw[arr] (sys.205) -- node[below=6pt, sloped, align=center, pos=0.5]{%
  \small Emotion Label (POS/NEG/NEU)\\[1pt]
  \small or AMBIGUOUS Result} (user.330);

\end{tikzpicture}
\caption{Level 0 Context Diagram}
\label{fig:dfd0}
\end{figure}

As shown in \textbf{Figure~\ref{fig:dfd0}}, the system acts as a black box that accepts raw textual data
and returns a categorized emotional state. The inclusion of an ``Ambiguous'' label differentiates
this system from standard classifiers, as it explicitly communicates uncertainty back to the user.

\subsection{Level 1: Detailed Data Flow}

The Level 1 DFD breaks down the system into its primary functional processes and shows the
internal data stores accessed during the inference phase.

%% ---- Figure 4.5: Level 1 DFD ----
\begin{figure}[H]
\centering
\begin{tikzpicture}[
  font=\small,
  box/.style={draw, rectangle, very thick, align=center,
              minimum width=2.8cm, minimum height=1.6cm, fill=green!10},
  datastore/.style={draw, rectangle, thick, align=center,
                    minimum width=2.8cm, minimum height=1.2cm, fill=purple!10},
  proc/.style={draw, circle, very thick, align=center, minimum size=3.0cm, fill=blue!8},
  arr/.style={-{Stealth[length=6pt, width=4pt]}, thick},
  lbl/.style={font=\scriptsize, midway}
]

% Nodes - Increased spacing to extend arrows
\node[box] (user)   at (0, 0)    {\textbf{User}\\\textbf{Input}};
\node[proc]         (feat) at (5, 0)    {\textbf{1.0}\\[4pt]\textbf{Feature}\\[2pt]\textbf{Extraction}};
\node[proc]         (emot) at (10.5, 0)  {\textbf{2.0}\\[4pt]\textbf{Emotion}\\[2pt]\textbf{Prediction}};
\node[datastore]    (model) at (5,-4.5) {\textbf{D1: Saved Model}\\[2pt](.joblib)};
\node[proc]         (ambi) at (10.5,-4.5) {\textbf{3.0}\\[4pt]\textbf{Ambiguity}\\[2pt]\textbf{Analysis}};
\node[box, fill=orange!12]  (out)  at (0,-4.5)  {\textbf{Final Output}\\\textbf{Label}};

% Arrows with adjusted label positions to avoid intersection
\draw[arr] (user.east)  -- node[lbl, above=3pt]{Text} (feat.west);
\draw[arr] (feat.east)  -- node[lbl, above=3pt]{Combined Features} (emot.west);
\draw[arr] (emot.south) -- node[lbl, right=3pt, align=left]{Probabilities\\\& Loop Counts} (ambi.north);
\draw[arr] (model.north) -- node[lbl, right=3pt, align=left]{Trained\\Params} (feat.south);
\draw[arr] (model.east)  -- node[lbl, above=3pt, align=center]{Classifiers\\Reference Clusters} (ambi.west);

% Final label path adjusted
\draw[arr] (ambi.south) -- ++(0,-0.8) -| node[lbl, below=3pt, pos=0.25, align=center]{Final Label} (out.south);

\end{tikzpicture}
\caption{Level 1 Data Flow Diagram (DFD)}
\label{fig:dfd1}
\end{figure}

\textbf{Figure~\ref{fig:dfd1}} details the internal flow of data. The process begins with
\textbf{Feature Extraction (1.0)}, which transforms raw text into a multi-modal feature vector
using parameters stored in \textbf{D1}. These features are consumed by \textbf{Emotion
Prediction (2.0)}, where XGBoost models generate class probabilities. Finally,
\textbf{Ambiguity Analysis (3.0)} uses both these probabilities and the topological loop counts
to decide whether to output a standard emotion label or an ambiguity flag.

\subsection{Level 2: Ambiguity Analysis Sub-Flow}

To further illustrate the core innovation of this project, we break down Process 3.0 into its
logical components. This Level 2 DFD shows how topological signals and classifier
probabilities are integrated to detect ambiguous inputs.

%% ---- Figure 4.6: Level 2 DFD ----
\begin{figure}[H]
\centering
\begin{tikzpicture}[
  font=\small,
  box/.style={
    draw, rectangle, very thick, align=center,
    minimum width=2.8cm, minimum height=1.1cm, fill=green!10
  },
  proc/.style={
    draw, circle, very thick, align=center,
    minimum size=3.0cm, text width=2.2cm, fill=blue!8
  },
  arr/.style={-{Stealth[length=6pt, width=4pt]}, thick}
]

%-----------------------------------------------
%  NODES  –  compact layout matching reference
%-----------------------------------------------

% Left input boxes
\node[box] (loop) at (0,  2.2) {Loop Counts ($H_1$)};
\node[box] (prob) at (0, -2.2) {Class Probabilities};

% Middle process circles
\node[proc] (ext)  at (5.5,  2.2) {%
  \textbf{3.1}\\[3pt]
  \textbf{Extreme}\\
  \textbf{Topology}\\
  \textbf{Check}
};
\node[proc] (conf) at (5.5, -2.2) {%
  \textbf{3.2}\\[3pt]
  \textbf{Confidence}\\
  \textbf{Gap}\\
  \textbf{Evaluation}
};

% 3.3 centred between 3.1 and 3.2
\node[proc] (asgn) at (10, 0) {%
  \textbf{3.3}\\[3pt]
  \textbf{Final Label}\\
  \textbf{Assignment}
};

% Result box
\node[box, fill=orange!12] (res) at (14, 0) {\textbf{Result}};

%-----------------------------------------------
%  ARROWS  –  labels well clear of every line
%-----------------------------------------------

% 1. Loop Counts --> 3.1 (Topology Counts) — straight horizontal
\draw[arr] (loop.east) -- (ext.west);
\node[font=\scriptsize, above=10pt]
  at ($(loop.east)!0.5!(ext.west)$) {Topology Counts};

% 2. Loop Counts --> 3.2 (Topology Counts) — vertical branch
%    Route: right from loop.east, drop down, then into conf at upper-left (150°)
%    This keeps it fully separate from the Probabilities arrow entering conf.west
\coordinate (Av) at ($(loop.east) + (0.8, 0)$);
\coordinate (Bv) at (Av |- conf.150);

\draw[arr] (loop.east) -- (Av) -- (Bv) -- (conf.150);

% Label to the LEFT of the vertical segment — large offset
\node[font=\scriptsize, left=10pt]
  at ($(Av)!0.5!(Bv)$) {Topology Counts};

% 3. Class Probabilities --> 3.2 (Probabilities) — straight horizontal
\draw[arr] (prob.east) -- (conf.west);
\node[font=\scriptsize, below=10pt]
  at ($(prob.east)!0.5!(conf.west)$) {Probabilities};

% 4. 3.1 --> 3.3 (Is Extreme?) — diagonal
\draw[arr] (ext.east) -- (asgn.north west);
\node[font=\scriptsize, above=12pt]
  at ($(ext.east)!0.5!(asgn.north west)$) {Is Extreme?};

% 5. 3.2 --> 3.3 (Confidence Gap) — diagonal
\draw[arr] (conf.east) -- (asgn.south west);
\node[font=\scriptsize, below=12pt]
  at ($(conf.east)!0.6!(asgn.south west)$) {Confidence Gap};

% 6. 3.3 --> Result — straight horizontal
\draw[arr] (asgn.east) -- (res.west);

\end{tikzpicture}
\caption{Level 2 DFD: Process 3.0 (Ambiguity Analysis)}
\label{fig:dfd2}
\end{figure}

The Level 2 DFD (\textbf{Figure~\ref{fig:dfd2}}) illustrates the internal decision logic of the ambiguity
detection system. It shows how H1 loop counts flow into two parallel processes: Process 3.1
handles extreme cases, while Process 3.2 evaluates the confidence gap between the primary and
secondary emotion classes. The outputs are reconciled in Process 3.3 to assign the final output
label.

\section{Training vs Inference}

Topology is computed during inference rather than precomputed. This ensures that the ambiguity
signal reflects the real structure of the training data relative to each input.
